\documentclass[11pt,letterpaper]{article}
\usepackage[utf8]{inputenc}
\usepackage[spanish,es-nodecimaldot]{babel}
\usepackage{geometry}
\geometry{top=2.5cm, bottom=2.5cm, left=2.8cm, right=2.8cm}
\usepackage{fontspec}
\usepackage{xcolor}
\usepackage{amsmath,amssymb}
\usepackage{booktabs}
\usepackage{array}
\usepackage{graphicx}
\usepackage{listings}
\usepackage{hyperref}
\usepackage{fancyhdr}

% Tipografías oficiales UD
\setmainfont{times.ttf}[
  Path = assets/fonts/,
  BoldFont = timesbd.ttf,
  ItalicFont = timesi.ttf,
  BoldItalicFont = timesbi.ttf
]
\newfontfamily\cambriafont{cambria.ttc}[
  Path = assets/fonts/,
  BoldFont = cambriab.ttf,
  ItalicFont = cambriai.ttf,
  BoldItalicFont = cambriaz.ttf
]

\definecolor{UDRed}{HTML}{ED3237}
\definecolor{UDDarkRed}{HTML}{B71C1C}
\definecolor{UDBlue}{HTML}{1291C0}
\definecolor{UDDarkGreen}{HTML}{047857}
\definecolor{UDSable}{HTML}{1E293B}
\definecolor{UDLightGray}{HTML}{F8FAFC}
\definecolor{UDBorder}{HTML}{CBD5E1}

\pagestyle{fancy}
\fancyhf{}
\fancyhead[L]{\small\color{UDSable}\cambriafont Maestría en Ciencias de la Información y las Comunicaciones (MCIC)}
\fancyhead[R]{\small\color{UDRed}\cambriafont Reto 1: Tubería CNN OpenACC}
\fancyfoot[C]{\thepage}
\renewcommand{\headrulewidth}{0.6pt}

\hypersetup{
  colorlinks=true,
  linkcolor=UDDarkRed,
  citecolor=UDBlue,
  urlcolor=UDBlue
}

\lstset{
  basicstyle=\ttfamily\small,
  backgroundcolor=\color{UDLightGray},
  frame=single,
  rulecolor=\color{UDBorder},
  breaklines=true,
  keywordstyle=\color{UDRed}\bfseries,
  commentstyle=\color{gray}\itshape,
  stringstyle=\color{UDDarkGreen},
  numbers=none,
  tabsize=2
}

\title{\vspace{-1.0cm}\Huge\bfseries\cambriafont\color{UDSable} Tubería CNN en OpenACC\\[0.3cm]
\Large\color{UDRed} Convolución 2D + ReLU + MaxPooling 2D\\[0.2cm]
\large\color{UDSable}\textmd{Documentación Técnica y Reporte de Rendimiento}}
\author{\textbf{\cambriafont Juferoga} \\
\textit{Computación Paralela y Distribuida} \\
Maestría en Ciencias de la Información y las Comunicaciones (MCIC) \\
Universidad Distrital Francisco José de Caldas}
\date{Septiembre de 2026}

\begin{document}
\maketitle
\thispagestyle{fancy}

\begin{abstract}
\noindent Este documento presenta el diseño, implementación y evaluación experimental de una tubería de operaciones fundamentales de Redes Neuronales Convolucionales (Convolución 2D de $100\times 100 \rightarrow 98\times 98$, activación ReLU y MaxPooling 2D con Stride 2 hacia $49\times 49$) acelerada en GPU mediante directivas OpenACC en C++. Se aborda la contención de datos en VRAM a lo largo de ráfagas iterativas ($N = 1\,000$ iteraciones) mediante la directiva \texttt{\#pragma acc data}, la maximización de la ocupación de los multiprocesadores CUDA mediante \texttt{collapse(2)}, la corrección teórica formal sobre la naturaleza topológica de la ventana de MaxPooling, y la evaluación cuantitativa frente a una línea base secuencial pura en CPU, alcanzando una aceleración de cómputo neto de \textbf{15.45x} y de sistema completo de \textbf{13.08x} con verificación numérica exacta ($\Delta_{\max} = 0.000$).
\end{abstract}

\vspace{0.4cm}
\tableofcontents
\vspace{0.6cm}
\hrule
\vspace{0.6cm}

\section{Introducción y Arquitectura de la Tubería}
En las capas convolucionales modernas, el procesamiento secuencial e iterativo de tensores espaciales demanda un uso intensivo de cómputo aritmético y ancho de banda de memoria. La secuencia estudiada en este reto consta de tres etapas fijas:
\begin{enumerate}
    \item \textbf{Convolución 2D:} Aplica un kernel espacial de $3 \times 3$ sobre una imagen de entrada $100 \times 100$ ($10\,000$ flotantes), generando una matriz de características de $98 \times 98$ ($9\,604$ flotantes).
    \item \textbf{Activación ReLU:} Transforma no linealmente cada elemento de la matriz anterior preservando su tamaño ($98 \times 98$), sustituyendo cualquier valor negativo por cero ($f(x) = \max(0, x)$).
    \item \textbf{MaxPooling 2D:} Realiza un submuestreo espacial disjunto mediante una ventana de $2 \times 2$ con paso (\textit{stride}) de 2, produciendo una salida final compacta de $49 \times 49$ ($2\,401$ flotantes).
\end{enumerate}

\section{Clarificación Conceptual y Corrección Teórica}
A diferencia del filtro convolucional (que sí posee una matriz de parámetros entrenables $3 \times 3$), la ``ventana de MaxPooling'' \textbf{no es una matriz de pesos ni requiere asignación de memoria dedicada en RAM/VRAM}. 
Corresponde estrictamente a un \textbf{rango de lectura e indexación espacial}:
\begin{equation}
\text{pool}(i, j) = \max_{(ki, kj) \in \{0, 1\}^2} \text{relu}(2i + ki, \; 2j + kj)
\end{equation}
La corrección teórica evita destinar memoria espuria o simular operaciones de convolución en el paso de submuestreo.

\section{Mapeo Canónico 1D (Row-Major Order)}
Para optimizar el uso de cachés y habilitar accesos coalescentes en la memoria global de la GPU, todas las matrices bidimensionales se asignan en memoria dinámica contigua (\texttt{new float[N]}):
\begin{itemize}
    \item \textbf{Imagen de Entrada ($100 \times 100$):} $\text{idx} = i \times 100 + j$, rango $[0 \dots 9\,999]$.
    \item \textbf{Buffer Convolución y ReLU ($98 \times 98$):} $\text{idx} = i \times 98 + j$, rango $[0 \dots 9\,603]$.
    \item \textbf{Salida MaxPooling ($49 \times 49$):} $\text{idx} = i \times 49 + j$, rango $[0 \dots 2\,400]$.
\end{itemize}

\noindent Para MaxPooling con Stride 2 ($W=98$), las 4 direcciones de lectura contigua son:
\begin{align*}
r_{00} &= (2i) \cdot 98 + 2j, & r_{01} &= (2i) \cdot 98 + (2j + 1) \\
r_{10} &= (2i + 1) \cdot 98 + 2j, & r_{11} &= (2i + 1) \cdot 98 + (2j + 1)
\end{align*}

\section{Estrategia de Paralelización en OpenACC}

\subsection{Ciclo de Vida de Datos y Persistencia en VRAM}
El bus PCIe impone una latencia crítica en sistemas acelerados. Para evitar transferencias redundantes en cada iteración, se emplea la directiva estructurada \texttt{\#pragma acc data}:
\begin{lstlisting}[language=C++]
#pragma acc data copyin(img[0:10000], kernel[0:9]) \
                 create(conv[0:9604], relu[0:9604]) \
                 copyout(pool[0:2401])
{
    for (int it = 0; it < iteraciones; ++it) {
        // Convolucion 2D (GPU) -> conv en VRAM
        // Activacion ReLU (GPU) -> relu en VRAM
        // MaxPooling 2D (GPU)  -> pool en VRAM
    }
}
\end{lstlisting}
\begin{itemize}
    \item \textbf{\texttt{copyin}:} Transfiere entrada y kernel a la GPU exactamente una vez al inicio.
    \item \textbf{\texttt{create}:} Aloja los buffers temporales \texttt{conv} y \texttt{relu} en VRAM con \textbf{cero tráfico PCIe}.
    \item \textbf{\texttt{copyout}:} Devuelve al Host únicamente la salida final compacta al terminar la ráfaga de 1,000 repeticiones.
\end{itemize}

\subsection{Colapso de Bucles Bidimensionales}
Cada etapa fusiona sus bucles anidados mediante \texttt{collapse(2)}, expandiendo el paralelismo para saturar los multiprocesadores CUDA:
\begin{itemize}
    \item Convolución 2D: $98 \times 98 = 9\,604$ hilos GPU.
    \item Activación ReLU: $98 \times 98 = 9\,604$ hilos GPU.
    \item MaxPooling 2D: $49 \times 49 = 2\,401$ hilos GPU.
\end{itemize}

\section{Resultados Experimentales y Ley de Amdahl}
A continuación se resumen las mediciones obtenidas para $N = 1\,000$ iteraciones:

\begin{table}[h!]
\centering
\cambriafont
\caption{Métricas de Rendimiento Experimental ($N = 1\,000$ iteraciones)}
\vspace{0.2cm}
\begin{tabular}{llr}
\toprule
\textbf{Métrica de Evaluación} & \textbf{Valor Medido} & \textbf{Unidad} \\
\midrule
CPU Secuencial ($T_{\text{CPU}}$) & $32.46$ & ms \\
GPU Cómputo Neto ($T_{\text{GPU\_Kernel}}$) & $2.10$ & ms \\
GPU Total con Transferencias ($T_{\text{GPU\_Total}}$) & $2.48$ & ms \\
Overhead de Bus PCIe & $0.38$ & ms ($15.3\%$) \\
\midrule
\textbf{Speedup de Cómputo ($S_{\text{cómputo}}$)} & \textbf{15.45x} & Aceleración \\
\textbf{Speedup Total de Sistema ($S_{\text{total}}$)} & \textbf{13.08x} & Aceleración \\
\midrule
\textbf{Consistencia Numérica ($\Delta_{\max}$)} & \textbf{0.000} & $0$ Discrepancias ($2\,401/2\,401$) \\
\bottomrule
\end{tabular}
\end{table}

\noindent \textbf{Análisis de Amdahl:} Si se ejecutase una sola iteración ($N = 1$), el costo de transferencias PCIe ($0.38$ ms) superaría con creces los microsegundos de cómputo neto, arrojando una desaceleración ($S < 1\times$). La ejecución en ráfaga iterativa amortiza el costo de transferencia a $0.38\ \mu\text{s}$ por iteración, posibilitando una aceleración real de más de \textbf{13x}.

\section{Conclusiones}
\begin{enumerate}
    \item OpenACC permite acelerar algoritmos de visión artificial y redes neuronales manteniendo una base de código compacta y legible en C++ estándar.
    \item La gestión explícita de datos con \texttt{\#pragma acc data} y la cláusula \texttt{create} es el pilar determinante de la eficiencia en arquitecturas heterogéneas.
    \item La verificación numérica bit a bit valida que la paralelización masiva de las etapas preserva de forma estricta la equivalencia funcional del modelo secuencial.
\end{enumerate}

\end{document}
